% Quarantine-only proof.
% Date: 29 August 2026.
% This file is not a live-book input and makes no novelty or priority claim.

\begin{proof}[Proof of
\cref{thm:ch14-cld-support-corrected-packet}]
We first prove
\cref{itm:ch14-cld-support-corrected-nontrivial}.  This is
\cite[Theorem~16.8.5]{CLD25}; we give the argument because the proper
tropical pushforward is essential and because it fixes the exponent without
any appeal to a ball-growth estimate on $X$.

The assertion is local.  Fix $x\in X$.  After shrinking, we may suppose
that $X$ is paracompact and that $\alpha$ is tropical.  We may also reduce to
the case in which $\alpha$ is positive.  Indeed, $T$ is symmetric, so its
wedge with the antisymmetric part of $\alpha$ is zero.  On a relatively
compact tropical chart, the symmetric form $\alpha$ is a difference of two
positive forms: if $\Theta$ is a smooth strictly positive $(p,p)$-form on a
neighborhood of the relevant compact set, then
$C\Theta+\alpha$ and $C\Theta$ are positive for $C\gg1$.  A finite smooth
partition of unity gives the same decomposition on an arbitrary compact
subset.  It is therefore enough to prove the assertion for each of the two
positive measures separately.

Choose a compact analytic neighborhood $V$ of $x$ and a moment map
\[
                         f:V\longrightarrow\mathbf G_m^N
\]
such that
\begin{equation}
\label{eq:ch14-cld-support-boundary-separation}
                 f_{\trop}(x)\notin f_{\trop}(\partial V)
\end{equation}
and such that $f$ tropicalizes $\alpha|_V$.  The separating moment exists
by \cite[Lemma~12.3.6]{CLD25}; concatenating it with a moment tropicalizing
$\alpha$ preserves \eqref{eq:ch14-cld-support-boundary-separation}.  Thus
there is a smooth $(p,p)$-superform $\omega$, defined near
$f_{\trop}(V)$, for which
\begin{equation}
\label{eq:ch14-cld-support-alpha-tropical}
                         \alpha|_V=f_{\trop}^*\omega.
\end{equation}
Multiplication by a plateau function permits us to regard $\omega$ as
compactly supported on $\mathbf R^N$ without changing it near the compact
set used below.

Set
\begin{equation}
\label{eq:ch14-cld-support-omega-U}
 \Omega=\mathbf R^N\setminus f_{\trop}(\partial V),
 \qquad
 U=f_{\trop}^{-1}(\Omega)\subset V.
\end{equation}
The compactness of $\partial V$ makes $\Omega$ open.  The set $U$ is
disjoint from $\partial V$, hence lies in $\operatorname{Int}(V/X)$ and is
open in $X$; by \eqref{eq:ch14-cld-support-boundary-separation}, it contains
$x$.  Moreover,
\begin{equation}
\label{eq:ch14-cld-support-proper-tropicalization}
                     f_{\trop}|_U:U\longrightarrow\Omega
\end{equation}
is proper: it is obtained from the proper map from the compact space $V$ to
$\mathbf R^N$ by restriction over $\Omega$.

Properness makes the following tropical pushforward well-defined:
\begin{equation}
\label{eq:ch14-cld-support-pushed-current}
 \langle T_\Omega,\beta\rangle
   :=\langle T|_U,f_{\trop}^*\beta\rangle,
 \qquad \beta\in A_c^{p,p}(\Omega).
\end{equation}
Pullback preserves the relevant positivity cone and commutes with
$\dprime$ and $\ddprime$.  Hence $T_\Omega$ is a positive closed
supercurrent of bidimension $(p,p)$ on $\Omega$.  In Lagerberg's bidegree
notation it has bidegree $(N-p,N-p)$.  Lagerberg's Euclidean support
theorem, \cite[Corollary~4.1]{Lag12Super}, says that the coefficient measures
of such a current vanish on every set of zero $p$-dimensional Hausdorff
measure.  (The same result is Corollary~3.5 in the arXiv version.)

The set
\[
                       E=f_{\trop}(A\cap U)
\]
has zero $p$-dimensional Hausdorff measure: it is contained in
$f_{\trop}(A\cap V)$, which is null by hypothesis.  It follows that all
coefficient measures of $T_\Omega$ vanish on $E$, and therefore the Radon
measure $T_\Omega\wedge\omega$ does not charge $E$.

For completeness, we now transfer this coefficientwise statement back to
$U$ without a cancellation argument.  We are treating the case in which
$\alpha$ is positive, so
\[
                           \mu=T\wedge\alpha
\]
is a positive Radon measure.  By
\eqref{eq:ch14-cld-support-alpha-tropical} and the definition
\eqref{eq:ch14-cld-support-pushed-current}, its ordinary measure-theoretic
pushforward is
\begin{equation}
\label{eq:ch14-cld-support-measure-pushforward}
                  (f_{\trop}|_U)_*(\mu|_U)
                         =T_\Omega\wedge\omega.
\end{equation}
Indeed, both sides take a compactly supported smooth function $h$ on
$\Omega$ to
$\langle T|_U,f_{\trop}^*(h\omega)\rangle$; equality for Radon measures
then follows from density of smooth compactly supported functions in
continuous compactly supported functions.

Let $B\subset A\cap U$ be compact.  Its image $C=f_{\trop}(B)$ is compact
and contained in $E$.  From
\eqref{eq:ch14-cld-support-measure-pushforward} we obtain
\[
 0=(T_\Omega\wedge\omega)(C)
   =\mu\bigl(f_{\trop}^{-1}(C)\cap U\bigr)
   \geq\mu(B).
\]
Thus $\mu(B)=0$.  Inner regularity gives $\mu(A\cap U)=0$.  Repeating this
argument for the two positive terms in the decomposition of $\alpha$ proves
\eqref{eq:ch14-cld-support-corrected-nontrivial} near $x$, and locality
proves it on $X$.

We next prove
\cref{itm:ch14-cld-support-corrected-trivial}.  Put
\begin{equation}
\label{eq:ch14-cld-support-p-equals-n-minus-q}
                              p=n-q.
\end{equation}
By multilinearity of the Bedford--Taylor product, it is enough to treat the
case in which every $u_j$ is psh.  The current $T$ is then positive and
closed.  The dpsh case follows by writing each $u_j$ locally as a difference
of two continuous psh functions and expanding; all support, pushforward,
and null-set assertions below are stable under finite linear combinations.

Fix $0<r<1$ and form the one-radius Gauss extension
\[
              K_r=K((\varpi)),\qquad |\varpi|=r.
\]
Equivalently, $K_r=\mathscr H(\eta_r)$ for the radius-$r$ Gauss point.  The
radius is $K$-free because $K$ is trivially valued and $r\ne1$.  In
particular, $K_r$ is nontrivially valued and universally multiplicative;
see \cite[Example~11.1.3 and \S12.7.4]{CLD25}.  Let
\[
 \pi:X_r=X\widehat\otimes_KK_r\longrightarrow X,
 \qquad
 \sigma:X\longrightarrow X_r
\]
be the projection and its continuous Shilov section.  Pure dimension and
boundarylessness are preserved by this scalar extension
\cite[\S\S11.2.5 and 11.4.1]{CLD25}, so the nontrivially valued part already
proved applies to $X_r$.

We need the following exact consequence of the Gauss construction:
\begin{equation}
\label{eq:ch14-cld-support-shilov-null}
        A\text{ tropically }H^p\text{-negligible}
        \quad\Longrightarrow\quad
        \sigma(A)\text{ tropically }H^p\text{-negligible in }X_r.
\end{equation}
Here is the verification.  The issue is local, so we may replace $X$ by an
affinoid domain.  Then $A$ is compact and $\sigma(A)$ is compact, hence
closed.  Let $W$ be a compact analytic domain of $X_r$ and
$f=(f_1,\ldots,f_m)$ a moment on $W$.  Applied to each $f_i$, and then
refined simultaneously, \cite[Lemma~16.8.7]{CLD25} gives a finite compact
cover $(C_\lambda)$ of $\sigma^{-1}(W)$, base moments $g_\lambda$ defined on
neighborhoods of $C_\lambda$, and vectors $e_\lambda\in\mathbf Z^m$ such
that
\begin{equation}
\label{eq:ch14-cld-support-gauss-moment-comparison}
 f_{\trop}\circ\sigma
       =g_{\lambda,\trop}+e_\lambda\log r
                 \quad\text{on }C_\lambda.
\end{equation}
After covering each $C_\lambda$ by finitely many compact analytic domains
on which $g_\lambda$ is a tropical chart, the hypothesis on $A$ shows that
$g_{\lambda,\trop}(A\cap C_\lambda)$ is $H^p$-null.  Translation by the
constant vector $e_\lambda\log r$ preserves Hausdorff measure.  The finite
union furnished by
\eqref{eq:ch14-cld-support-gauss-moment-comparison} proves
\eqref{eq:ch14-cld-support-shilov-null}.

Write
\[
 u_{j,r}=\pi^*u_j,
 \qquad
 T_r=\bigwedge_{j=1}^q\dprimeddprime u_{j,r},
 \qquad
 \alpha_r=\pi^*\alpha.
\]
The current $T_r$ has bidegree $(q,q)$ and therefore bidimension
$(p,p)$ by \eqref{eq:ch14-cld-support-p-equals-n-minus-q}.  Applying
\cref{itm:ch14-cld-support-corrected-nontrivial} on $X_r$ to the
$H^p$-null set $\sigma(A)$ gives
\begin{equation}
\label{eq:ch14-cld-support-upstairs-null}
                         |T_r\wedge\alpha_r|(\sigma(A))=0.
\end{equation}

Two exact scalar-extension identities now finish the proof.  Since
$\alpha_r$ has bidegree $(p,p)=(n-q,n-q)$,
\cite[Proposition~16.8.2]{CLD25} gives
\begin{equation}
\label{eq:ch14-cld-support-shilov-support}
                \operatorname{supp}(T_r\wedge\alpha_r)
                                  \subseteq\sigma(X).
\end{equation}
The pushforward formula in \cite[Proposition~8.7.5]{CLD25}, in the scalar
extension form recalled immediately after
\cite[Proposition~16.8.2]{CLD25}, gives
\begin{equation}
\label{eq:ch14-cld-support-gauss-pushforward}
                  \pi_*(T_r\wedge\alpha_r)=T\wedge\alpha.
\end{equation}
Because $\pi\circ\sigma=\operatorname{id}_X$, one has
\[
        \pi^{-1}(A)\cap\sigma(X)=\sigma(A).
\]
Together with
\eqref{eq:ch14-cld-support-upstairs-null} and
\eqref{eq:ch14-cld-support-shilov-support}, this says that
$T_r\wedge\alpha_r$ does not charge $\pi^{-1}(A)$.  Pushing forward through
\eqref{eq:ch14-cld-support-gauss-pushforward} proves
\eqref{eq:ch14-cld-support-corrected-trivial}.  Notice that the Hausdorff
exponent used upstairs is
$p=n-q$, the bidimension of $T_r$, and not the number $q$ of curvature
factors.
\end{proof}

\begin{proof}[Proof of
\cref{prop:ch14-cld-support-printed-exponent-counterexample}]
For $t\in\mathbf R$, let $s(t)$ be the Gauss seminorm
\[
 \left|\sum_{j\in\mathbf Z}a_jz^j\right|_{s(t)}
       =\max_{a_j\ne0}\mathrm e^{jt}.
\]
Its image is the canonical skeleton of $X$, and
$\tau(s(t))=t$; in particular $s(0)=\eta_1$.  The function
$h(t)=\max\{0,t\}$ is continuous and convex.  Hence
$u=h\circ\tau$ is CLD-psh by the tropical convexity criterion
\cite[Proposition~8.2.5]{CLD25}.

We compute its curvature with the lattice coordinate $t$.  The
one-dimensional tropical current formula identifies the curvature of this
tropical potential with the distributional Hessian of $h$ on the canonical
skeleton.  Explicitly, if $\varphi$ is a compactly supported smooth test
function, the calibrated integration formula and two integrations by parts
give
\begin{align}
 \left\langle\dprimeddprime[u],\varphi\right\rangle
   &=\int_{\mathbf R}\max\{0,t\}
           \frac{\mathrm d^2}{\mathrm dt^2}
             \bigl(\varphi\circ s\bigr)(t)\,\mathrm dt \notag\\
   &=-\int_0^\infty
           \frac{\mathrm d}{\mathrm dt}
             \bigl(\varphi\circ s\bigr)(t)\,\mathrm dt
     =\varphi(s(0)).
\label{eq:ch14-cld-support-kink-test}
\end{align}
There is no boundary term because $\varphi$ has compact support.  Equivalently,
this is the slope jump $1-0=1$ at the only corner of $h$.  The tropical
Bedford--Taylor construction, or directly
\cite[Theorem~8.10.3]{CLD25}, shows that no curvature is present off this
zero-dimensional skeletal corner.  Thus
\eqref{eq:ch14-cld-support-kink-test} proves
\eqref{eq:ch14-cld-support-kink-dirac}.

Every tropical chart sends the singleton $A=\{\eta_1\}$ either to the empty
set or to a singleton.  Its one-dimensional Hausdorff measure is therefore
zero, so $A$ is tropically $H^1$-negligible.  On the other hand,
\[
                    \dprimeddprime[u](A)=1.
\]
For $n=q=1$ and $\alpha=1$, the exponent printed in
\cite[Corollary~16.8.6]{CLD25} is $H^q=H^1$ and would force this mass to be
zero.  This contradiction proves that the printed exponent cannot be
correct.  With the corrected exponent $H^{n-q}=H^0$, there is no
contradiction: a nonempty singleton has positive $H^0$-measure.
\end{proof}
